\chapter{Conclusiones y recomendaciones}


La presente investigación tuvo como objetivo estudiar la interdependencia productiva internacional y su relación con los procesos de integración comercial regional mediante herramientas de análisis de redes. Para ello, se construyó un conjunto de redes insumo--producto internacionales a partir de las tablas \gls{icio} de la \gls{ocde} y se aplicaron métodos de detección de comunidades con el propósito de identificar posibles bloques comerciales formados por sectores económicos que mantienen relaciones productivas más intensas entre sí que con el resto de la economía mundial. 

La importancia de este estudio deriva de los profundos cambios que ha experimentado la economía mundial durante las últimas décadas. La globalización ha incrementado la interdependencia económica entre los países y favorecido la fragmentación de la producción mundial, así como su articulación en cadenas globales de valor. Paralelamente, el comercio internacional se desarrolla en un contexto marcado por la tensión entre la creciente liberalización de los mercados a escala internacional y los diversos procesos de integración económica regional. Este contexto plantea la necesidad de analizar cómo se estructura el sistema productivo internacional, más allá de los acuerdos comerciales y económicos formales.

Las matrices insumo--producto multiregionales (\glspl{mriot}) ofrecen una herramienta útil para analizar la estructura productiva mundial, ya que integran en un único marco contable los flujos intersectoriales nacionales e internacionales. Dentro de estos instrumentos, la base de datos \gls{icio}--\gls{ocde} destaca por su amplia cobertura sectorial, geográfica y temporal, al incorporar información de 
80 países y 50 sectores económicos para el periodo de 1995 a 2022. Además, asegura la consistencia metodológica y la armonización de los datos estimados, lo que la convierte en una fuente adecuada para el estudio de la interdependencia productiva internacional.

Para realizar el análisis, los datos de la submatriz de consumo intermedio de las \glspl{mriot} de \gls{icio}--\gls{ocde} se representaron por medio de redes dirigidas y ponderadas. En estas redes insumo--producto internacionales, cada vértice representa un sector económico de un país, mientras que cada arco representa una relación productiva dirigida desde un sector suministra bienes o servicios intermedios hacia otro sector que los consume productivamente, ya sea dentro de una economía nacional o a través del comercio internacional. Por su parte, el peso asociado a cada arco corresponde al valor del flujo comercial normalizado respecto al valor total de las transacciones intersectoriales registradas en el año respectivo.

La detección de comunidades es un enfoque de la teoría de redes que busca particionar los vértices de una red en comunidades, de tal forma que los vértices de una comunidad tengan mayor densidad de conexiones entre sí, que con los vértices de otras comunidades. Esta vinculación puede medirse en términos de las aristas de la red o de la intensidad de sus pesos. Uno de los criterios más utilizados en la detección de comunidades es la maximización de la modularidad, una medida que compara la intensidad de las conexiones observadas dentro de las comunidades con las conexiones que se esperarían si los vínculos de la red se distribuyeran aleatoriamente sin formar comunidades. Entre los métodos de detección de comunidades que maximizan la modularidad destacan los algoritmos de Leiden y de Louvain, por su eficiencia computacional y la calidad de sus soluciones. Ambos algoritmos se basan en una estrategia multinivel que alterna etapas de búsqueda local y agregación de comunidades con el objetivo de mejorar iterativamente la partición obtenidas.

En la literatura existen diversos estudios que han aplicado métodos de detección de comunidades al análisis de redes de comercio internacional, incluidas las redes insumo--producto, con el propósito de identificar posibles bloques comerciales. En comparación con esos trabajos, esta tesina amplía la cobertura sectorial, geográfica y temporal al utilizar las tablas \gls{icio}--\gls{ocde} versión 2025. Además, incorpora el algoritmo de Leiden al análisis de redes insumo--producto internacionales, uno de los métodos más avanzados de detección de comunidades.

El primer hallazgo de esta investigación es la existencia de una fuerte estructura comunitaria en las redes insumo-producto internacionales durante el periodo 1995--2022.  Los valores de modularidad dirigida y ponderada obtenidos mediante el algoritmo de Leiden son altos en comparación a los rangos reportados en la literatura para redes con estructura comunitaria bien definida. Esto indica que la economía mundial se encuentra estructurada en bloques productivos con elevados niveles de interdependencia interna. Las comunidades o bloques identificados corresponden a economías nacionales individuales o a la unión de varias economías nacionales, más que a cadenas productivas sectoriales organizadas a nivel internacional. Estos resultados coinciden con hallazgos previos reportados en estudios que utilizan otras bases de datos.

\textcolor{red}{Estos resultados son relevantes porque sugieren  que, a pesar de la creciente fragmentación internacional de la producción y la consolidación de las \glspl{cgv} asociadas a determinadas industrias, como la automotriz o la electrónica, la interdependencia productiva capturada por las redes insumo--producto internacional se estructura en términos nacionales y regionales. Desde esta perspectiva, los métodos de detección de comunidades permiten identificar aquellos casos en los que esta interdependencia trasciende las fronteras nacionales y da lugar a bloques regionales caracterizados por elevados niveles de integración intersectorial. Lo anterior no contradice la existencia de las \glspl{cgv}, sino que sugiere que estos vínculos se encuentran integrados dentro de espacios productivos nacionales o regionales.}

\textcolor{red}{Debido a que las comunidades detectadas se componen de una o varias economías nacionales, este enfoque permite analizar la integración productiva regional en dos sentidos. En primer lugar, permite evaluar la evolución y estabilidad de los bloques productivos en términos de las economías nacionales que los integran anualmente; así como comparar estos patrones con el desarrollo de los acuerdos formales de integración económica.   Por ejemplo, el estudio de caso del sudeste asiático mostró un primer periodo de consolidación regional de 1995 a 2008, seguido por una etapa de mayor inestabilidad que derivó en la fragmentación de la región en tres comunidades estables: una conformada por Corea del Sur, Vietnam y Camboya, una comunidad individual correspondiente a Indonesia y otra integrada por el resto de los países de la \gls{asean}. En contraste, en el caso de América del Norte, los resultados muestran que, a pesar de la existencia del \gls{tlcan} y posteriormente del \gls{tmec}, las tres economías permanecen como comunidades independientes durante todo el periodo analizado. Esto muestra que los procesos de integración económica no necesariamente se traducen en bloques productivos plenamente cohesionados. En segundo lugar, este enfoque permite identificar sectores específicos que se integran a comunidades distintas de aquella a la que pertenece la mayoría de los sectores de su país. Estos casos pueden interpretarse como indicios de la integración de ciertos sectores en \glspl{cgv} o en espacios productivos externos a su economía nacional.}

A nivel metodológico, los resultados muestran que el algoritmo de Leiden obtuvo mejores resultados que el algoritmo de Louvain, utilizado en otros estudios similares dentro de la literatura. Además, las soluciones producidas por el algoritmo de Leiden mostraron gran robustez frente al tipo de partición inicial utilizado y una menor variabilidad en los tiempos de ejecución. Esto confirma la pertenencia de incorporar el algoritmo de Leiden al análisis de redes insumo--producto.

\textcolor{red}{En conjunto, este estudio muestra que la detección de comunidades en redes insumo--producto internacionales constituye una herramienta útil para analizar la interdependencia productiva y su relación con los procesos de integración regional. Las comunidades identificadas permiten aproximar la existencia de bloques productivos con fuertes vínculos de integración comercial, que no necesariamente coinciden con los acuerdos comerciales formales, lo que permite diferenciar la integración institucional establecida por medio de acuerdos comerciales y la integración observada a partir de las relaciones productivas efectivas.}

% Limitaciones 
% Lineas futuras de investigación 